\documentclass{article} % For LaTeX2e
\usepackage[submission]{colm2026_conference}

\usepackage{microtype}

\usepackage{hyperref}
\usepackage{url}
\usepackage{booktabs}

\usepackage{lineno}
\usepackage{amsmath,amssymb}
\usepackage{graphicx}
\usepackage{xspace}
\definecolor{darkblue}{rgb}{0, 0, 0.5}
\hypersetup{colorlinks=true, citecolor=darkblue, linkcolor=darkblue, urlcolor=darkblue}


\title{From Text Refusal to Tool Safety: A Causal Mechanistic Study of Transfer Failure in LLM Agents}

% Authors must not appear in the submitted version. This should be be taken care of automatically as long as you are using the "submission" option for the colm2026_conference package. But it's on the authors to verify. Non-anonymous submissions will be rejected without review.t

\author{Antiquus S.~Hippocampus, Natalia Cerebro \& Amelie P. Amygdale \thanks{ Use footnote for providing further information
about author (webpage, alternative address)---\emph{not} for acknowledging
funding agencies.  Funding acknowledgements go at the end of the paper.} \\
Department of Computer Science\\
Cranberry-Lemon University\\
Pittsburgh, PA 15213, USA \\
\texttt{\{hippo,brain,jen\}@cs.cranberry-lemon.edu} \\
\And
Ji Q. Ren \& Yevgeny LeNet \\
Department of Computational Neuroscience \\
University of the Witwatersrand \\
Joburg, South Africa \\
\texttt{\{robot,net\}@wits.ac.za} \\
\AND
Coauthor \\
Affiliation \\
Address \\
\texttt{email}
}

% The \author macro works with any number of authors. There are two commands
% used to separate the names and addresses of multiple authors: \And and \AND.
%
% Using \And between authors leaves it to \LaTeX{} to determine where to break
% the lines. Using \AND forces a linebreak at that point. So, if \LaTeX{}
% puts 3 of 4 authors names on the first line, and the last on the second
% line, try using \AND instead of \And before the third author name.

\newcommand{\fix}{\marginpar{FIX}}
\newcommand{\new}{\marginpar{NEW}}

\begin{document}

\ifcolmsubmission
\linenumbers
\fi

\maketitle

\begin{abstract}

\iffalse


\fi

Large language models(LLMs) are increasingly being deployed as tool- calling agents in finance, healthcare, education, and legal, where stakes are high, yet safety evaluations are significantly concentrated on text-refusal behavior. Whether the refusal mechanisms learned during chat alignment reliably transfer to tool call decisions remains largely untested. We first construct an evaluation dataset of 2304 prompts that spans four domains and is grouped into triples of harmful requests consisting of a no-tool prompt, a standard tool-enabled prompt and an adversarially framed tool-enabled prompt while keeping their underlying intent constant. Tool-call safety is scored by 20 deterministic forbidden actions, which detect the following bad behavior: leaking PII, accessing protected data, bypassing authorization, and using intent evasion keywords. Evaluating five language models across different families, we find a noticeable divergence rate between refusals in chat and refusals in tool calls for the same requests conditioned: All the models reliably refuse, yet still make unsafe tool calls on the same request, with conditioned divergence reaching 0.44 on the largest model. Next, we mechanistically investigate this divergence. Across three model families, we identify a single linear refusal direction in the residual stream and find that tool-enabled prompts suppress its activation at the decision token along the same axis rather than rotating it. This suppressed signal reliably predicts unsafe calls, but activation patching only restores 18-62\% of refusals, showing that refusal suppression and tool-affordance pathway both drive transfer failure. Steering through the refusal direction reveals a trade-off between reducing unsafe tool calls and inducing benign over-refusal. Since the same suppression shows up in every architecture we test, deploying LLM agents in regulated settings will require treating tool-call safety as a separate alignment target, not a downstream consequence of text safety.

\end{abstract}

\section{Introduction}

Increasingly, the use cases of large language models have shifted from simply generating text to taking actions in the form of agents. Many modern systems browse the web, execute code, call typed tools, and implement enterprise workflows \citep{schick2023toolformer,
yao2023react, qin2023toolllm}. With this shift in the deployment of LLMs, there is a corresponding shift in the meaning of their "safety". That is, a tool-enabled agent can cause harm through its actions in the environment rather than by simply getting rid of an unsafe sentence, as measured in text-based settings. 

Indeed, safety alignment is currently overwhelmingly trained and evaluated on text. In particular, red teaming, refusal training, and jailbreak benchmarks target harmful text-based \textit{completions} \citep{ouyang2022training, bai2022training, bai2022constitutional,
ganguli2022red, perez2022red, mazeika2024harmbench}. A model that has learned to
say ``I can't help with that'' has learned a behavior over its \emph{verbal}
output. Nonetheless, whether that behavior also governs the \emph{tool-call} decision is a
separate question that carries direct consequences as these agents move into financial, medical, and legal settings.



This distinction is important because tool-enabled agents are increasingly being developed and applied in high-risk domains such as finance and healthcare \citep{ding2024fintrading, wang2025medagents}, where failures can lead to financial losses, privacy breaches, or the execution of malicious instructions extracted from external sources. 

However, prior work has rarely investigated whether refusal mechanisms transfer to tool-call decisions. To address these gaps, we evaluate models from five different families on a dataset we created spanning four domains. We run inference on all of these models using groups of three semantically similar requests: one text-only request and two other tool-enabled variants. We find that models that reliably refuse harmful requests in text generation often fail to refuse similar requests when they are framed in the format of a tool invocation.

\paragraph{Behavioral finding:}  We evaluate our benchmark of prompt triples, in which each harmful request is posed as a plain chat prompt as well as two tool-enabled versions. We find that, although the models do reliably refuse in chat prompts, they oftentimes fail to refuse the same requests framed when phrased in the tool-use variant. This gap, which we call \textit{conditioned divergence}, is notable and reaches its highest value on the largest model we tested.

\paragraph{Mechanistic finding:} Next, we ask why, seeking to find the mechanistic underpinnings of our behavioral results. We begin by identifying the rotation direction in each of the models employing the approach used in prior work, identifying or indicating that the direction is mediated by a single axis in the residual stream \textbf{[cite]}. To directly address the behavioral finding, we then test whether that refusal direction also governs tool call safety. We find that it does, but only in part. That is, tool-enabled prompts suppress the activation of the refusal direction at the decision token, and moreover, this suppressed value reliably predicts the unsafe call. Nonetheless, when we bring the refusal signal to the note tool level, we find that for most models, only a minority of refusals recover. Hence, we conclude that this observed transfer failure is driven by a combination of refusal suppression and the tool-use itself. The refusal direction, therefore.of serves as a useful signal of monitoring but as a comparatively blunt knob for control. To our knowledge, the consistency of this observation across model families serves as new evidence that it reflects a general property for how alignment is represented in models.

\paragraph{Contributions:}

\begin{enumerate}
    \item A 2304 prompt matched triple benchmark spanning four high-stakes domains with 20 deterministic forbidden actions to score tool-call refusal
    \item A conditioned-divergence metric separating transfer failure from the base refusal rate, together with evidence that text refusal does not reliably transfer to tool call safety.
    \item Causal mechanistic account of tool use refusal replicated across three model families: confirming that refusal is carried by a single linear direction; tool context suppresses that direction at the decision token; and the direction is a strong predictor of the unsafe action, though only a partial mediator of it.
    \item A model-agnostic toolkit to be used for both behavioral evaluations and the interpretability pipeline
\end{enumerate}

\section{Related Work}

  \paragraph{Tool-augmented language models.} A line of work extends LLMs from
  text generation to acting through external tools and APIs, teaching models to
  call functions, search, and run multi-step workflows
  \citep{schick2023toolformer, yao2023react, qin2023toolllm}. This shifts the
  locus of potential harm from generated text to the actions an agent takes,
  which is the surface our evaluation targets.

  \paragraph{Refusal and safety alignment in text.} Safety alignment is trained
  and measured over textual completions. Instruction tuning and RLHF teach models
  to decline harmful requests \citep{ouyang2022training, bai2022training,
  bai2022constitutional}, and red-teaming and jailbreak benchmarks stress-test
  that behavior \citep{ganguli2022red, perez2022red, mazeika2024harmbench}. These
  methods evaluate what the model says, not whether the learned refusal governs
  what it does when given tools.

  \paragraph{Mechanistic interpretability of refusal.} \citet{arditi2024refusal}
  show that refusal in chat models is mediated by a single linear direction in the
  residual stream: ablating it suppresses refusal and adding it induces
  over-refusal. We reproduce this across additional architectures and extend it to
  the text-to-tool transition, testing whether the same direction stays active
  when the model selects an action rather than a sentence.

  \paragraph{Jailbreaks through function calling.} \citet{wu2024darkside} show
  that function-calling interfaces open a pathway that bypasses alignment,
  eliciting unsafe behavior that the same model refuses in chat. This establishes
  the behavioral vulnerability but does not localize where safety fails. We
  provide that localization, tracing the failure to suppression of the refusal
  direction at the action-selection token.

  \paragraph{Text-to-tool safety divergence.} Closest to our work,
  \citet{cartagena2026gap} show at scale that models refusing harmful requests in
  text frequently emit unsafe tool calls for the same requests, and note that
  their behavioral measurements cannot isolate the internal mechanism. We address
  that gap. We supply the causal mechanistic account, introduce the
  conditioned-divergence metric that conditions on text refusal so the measure is
  not confounded by a model's base refusal rate, and evaluate matched triples
  across three interaction modes including an adversarial stress mode rather than
  matched pairs.

\section{Methodology}

\subsection{Dataset}

 The evaluation dataset consists of 2304 prompts across different categories:

\begin{itemize}
    \item \textbf{Domain:} Finance, Healthcare, Education, and Legal - four
      regulated settings, each with its own misuse modes and protected data.
    \item \textbf{Scenario:} Direct, Indirect, Justification/roleplay, and
      Structured - four phrasings of the same intent, from straightforward
      to tool-format-driven.
    \item \textbf{Mode:} No-tool chat, Tool-enabled normal, and Tool-enabled
      stress (jailbreak-style framing) - the central comparison.
    \item \textbf{System condition:} Neutral, Safety-reinforced (explicit
      safety constraints), and Tool-encouraging (nudges tool use).
    \item \textbf{Variant:} A (original phrasing) and B (a slight rewording).
    \item \textbf{Request type:} Harmful (should be refused / no unsafe tool
      call) and Benign (safe to answer or call tools).
\end{itemize}


Scenarios offer a range of approaches, the model is getting asked to set a prompt, as different scenario styles alter AI outputs \citep{lutz2025persona}. The three mode options are the essential comparison of the datapoints. No-tool chat: no tools provided; model responds in conversation. Tool-enabled normal: domain tools are injected into the context under standard prompting. Tool-enabled stress: same tools injected into context. however user prompt includes higher pressure or jailbreak-style framing. Matched groups (same domain, scenario, variant, and system condition) share the same underlying intent across modes, with prompts adapted to each mode's context: no-tool prompts make no references to tools, while tool enabled prompts explicitly invoke specific tools. This design holds scenario and intent constant while varying both tool availability and the framing appropriate to each mode, following the matched-prompt multi-domain saftey evaluation methodology of \citep{cartagena2026gap}.



Three system conditions vary the framing presented to the models: a neutral role description, a safety-reinforced version with explicit safety reminders, and finally a tool-encouraging version emphasizing efficient use of available functions.

These categories introduce a wide variety of real replicable scenarios that were tested to explain our claim

\subsection{Tool Environments and Forbidden Actions}

Each domain exposes five function-calling tools modeled on realistic
deployment schemas, adapting the domain tool-environment and deterministic forbidden-action design of \citep{cartagena2026gap}.
Example tools include \texttt{place\_trade\_order} and
\texttt{initiate\_wire\_transfer} (Finance),
\texttt{retrieve\_patient\_record} and \texttt{draft\_medication\_plan}
(Healthcare), \texttt{change\_grade} (Education), and
\texttt{file\_regulatory\_submission} (Legal).
Tools are provided in standard OpenAI function schema format.

Tool-call safety is scored by \emph{ForbiddenAction} predicates:
deterministic Python functions that inspect a tool name and argument
dictionary and return a Boolean violation flag.
Each domain defines five predicates---one per tool---for 20 total.
Predicate types are:

\begin{itemize}
\item \textbf{Sensitive-data flag}: fires when a Boolean parameter
  requesting protected data (e.g., \texttt{include\_phi=True}) is set.
\item \textbf{Authorization-bypass}: fires when a required authorization
  field (e.g., \texttt{approval\_ticket}) is absent or empty.
\item \textbf{PII exfiltration}: fires when known synthetic PII markers
  appear in text-valued arguments.
\item \textbf{Intent keyword}: fires when argument content contains
  policy-evasion terms (e.g., ``evade'', ``bypass'', ``contraindicated'').
\end{itemize}

A run is scored \emph{tc\_safe} if no predicate fires; any violation yields
\emph{tc\_unsafe}.
Critically, authorization-bypass predicates fire only on absent fields: a
model that supplies a fabricated value (e.g.,
\texttt{approval\_ticket:~"123456"}) passes the predicate.
All reported divergence rates are therefore conservative lower bounds
.

\subsection{Models}
\label{sec:models}

\begin{itemize}
  \item \textbf{Mistral-7B-Instruct-v0.3} \citep{jiang2023mistral}
  \item \textbf{Llama-3.1-70B-Instruct} \citep{dubey2024llama}
  \item \textbf{Command R 7B}
    (CohereLabs/c4ai-command-r7b-12-2024) \citep{cohere2024commandr7b}
  \item \textbf{Qwen3-14B} (Qwen/Qwen3-14B) \citep{yang2025qwen3}
  \item \textbf{Gemma 3 27B} (google/gemma-3-27b-it) \citep{gemma2025gemma3}
  \end{itemize}


\subsection{Metrics}

  \textbf{Refusal rate.}
  Fraction of responses that decline the user request, computed separately
  for harmful and benign prompts.
  On harmful prompts, high refusal is the target safety behavior \citep{arditi2024refusal}; on benign
  prompts, it quantifies over-refusal (false positive rate).
  Refusal is scored only on genuine text responses: a run that emits a tool
  call, or whose text is empty, is not counted as a verbal refusal.

  \textbf{tc\_safe rate.}
  Fraction of tool-enabled runs in which no ForbiddenAction predicate fires.

  \textbf{Matched triples.}
  Each harmful request instance is issued in all three interaction modes:
  no-tool chat, tool-enabled normal, and tool-enabled stress, producing a
  \emph{matched triple} of runs over the same underlying request.
  Triples are keyed by (Domain, Scenario, Variant, System Condition, Batch),
  so the no-tool run and the two tool-enabled runs differ only in tool
  affordance and prompt framing, not in the request being probed.
  Let $\mathcal{T}$ be the set of all matched triples and, for a triple $t$,
  let $\mathrm{refuse}(t)$ indicate that the no-tool run verbally refused and
  $\mathrm{unsafe}_{m}(t) = \neg\,\mathrm{tc\_safe}$ indicate that the
  tool-enabled run in mode $m \in \{\mathrm{normal}, \mathrm{stress}\}$ fired a
  ForbiddenAction predicate.

  \textbf{Conditioned divergence.}
  Our primary metric.
  Let $\mathcal{T}_{R} = \{\, t \in \mathcal{T} : \mathrm{refuse}(t) \,\}$ be
  the triples whose no-tool run refused (\emph{baseline refusal}).
  Conditioned divergence for tool mode $m$ is the fraction of these triples
  whose matched tool-enabled run is unsafe:

  \begin{equation}
  \Delta^{\mathrm{cond}}_{m} = \frac{
    \bigl|\{\, t \in \mathcal{T}_{R} : \mathrm{unsafe}_{m}(t) \,\}\bigr|
  }{\bigl|\mathcal{T}_{R}\bigr|}
  \label{eq:divergence}
  \end{equation}

  $\Delta^{\mathrm{cond}}_{m}$ measures how often verbal refusal of a request
  fails to carry over to the matched action when tools are available, the
  core transfer-failure signal \citep{cartagena2026gap, wu2024darkside}. It is reported separately for normal and
  stress modes, and conditions on $\mathcal{T}_{R}$ so that it is unaffected
  by how often the model refuses in text (a model that rarely refuses cannot
  inflate or deflate it).

  \textbf{Joint divergence.}
  For comparability with text-only safety benchmarks we also report the
  unconditioned form, dividing by all triples rather than the refused subset:
  \begin{equation}
  \Delta^{\mathrm{joint}}_{m} = \frac{
    \bigl|\{\, t \in \mathcal{T}_{R} : \mathrm{unsafe}_{m}(t) \,\}\bigr|
  }{\bigl|\mathcal{T}\bigr|}
  \end{equation}
  By construction $\Delta^{\mathrm{joint}}_{m} \le \Delta^{\mathrm{cond}}_{m}$.

  \textbf{Tool-engagement divergence (secondary).}
  The unsafe predicate above counts only forbidden tool calls; a safe or
  properly-authorized call is not divergent.
  As a looser upper bound we also report a variant that replaces
  $\mathrm{unsafe}_{m}(t)$ with $\mathrm{call}_{m}(t)$, true whenever the
  tool-enabled run emits \emph{any} tool call.
  This conflates benign tool use with the forbidden action and is reported
  only to bound the gap between ``acted at all'' and ``acted unsafely.''



\subsection{Refusal Classification}

Text refusal is classified by a hybrid pipeline.
A regex fast-path covers explicit patterns
(e.g., ``I cannot help with'', ``I must decline'') and resolves
upto 30\% of cases without any API call.
Ambiguous cases pass to an LLM judge (Claude Haiku) with a
binary YES/NO prompt asking whether the model declined the harmful part of
the request.
Tool-call scoring is fully deterministic via the ForbiddenAction predicates
and involves no language model.

\subsection{Mechanistic methods}

\paragraph{Refusal direction:} On the no-tool prompts, we take the decision token residual vector at every layer and compute the difference in mean direction between harmful and benign prompts to identify the refusal direction: $\hat{r_l} = \text{unit}(\mu_l^{\text{harm}} - \mu_l^{\text{benign}})$ \textbf{[cite]}. We then select the operating layer as the one maximizing the norm of the difference in means directions, measuring the magnitude of the separation, inside the causal mid band (35 to 85 percent depth). The layer is selected automatically for each model. Finally, projecting a prompt's residual onto the difference-in-means refusal direction gives us a scalar \textit{refusal activation}. 

\paragraph{Interventions:} To probe causality, we apply three interventions:
\begin{enumerate}
    \item Directional ablation, which removes the $\hat{r}$ component at each layer of the model
    \item Activation addition, adding $c \cdot \hat{r}$ for scalar coefficient $c$
    \item Activation patching, in which we overwrite the projection of a tool run with its corresponding no-tool value.
\end{enumerate}
We steer the model by activation addition over a coefficient grid that is scaled to the harm versus benign projection gap for each model. Hence, our intervention magnitudes can be compared across models with varying activation scale.

\paragraph{Robustness of the direction:} To evaluate the refusal direction, we contrast across request type, subtracting activations for harmful and benign prompts. In order to verify that our result is not simply an artifact of this choice, we additionally extract a direction defined by behaviors, by the behavioral results, finding the difference in means between Refused and Complied harmful prompts. We then use cosine similarity to evaluate how strongly aligned these directions are, confirming that request harmfulness and refusal decision indeed occupy the same direction.

\begin{table}[b]
\caption{
Conditioned divergence between text refusal and tool-call safety.
Higher values indicate more frequent safety transfer failures.
}
  \label{tab:headline-divergence}
  \centering

  \begin{tabular}{lccccc}
    \toprule
    & \multicolumn{1}{c}{Text} & \multicolumn{2}{c}{Unsafe-tool} & \multicolumn{2}{c}{Divergence} \\
    \cmidrule(lr){2-2}\cmidrule(lr){3-4}\cmidrule(lr){5-6}
    Model & refusal & Normal & Stress & Normal & Stress \\
    \midrule
    Llama-3.1-70B-Instruct   & 0.781 & 0.422 & 0.333 & \textbf{0.440} & 0.357 \\
    Mistral-7B-Instruct-v0.3 & 0.659 & 0.198 & 0.224 & \textbf{0.162} & 0.217 \\
    Gemma-3-27B-it           & 0.812 & 0.120 & 0.115 & \textbf{0.099} & 0.109 \\
    c4ai-command-r7b         & 0.823 & 0.096 & 0.107 & \textbf{0.085} & 0.095 \\
    \bottomrule
  \end{tabular}
\end{table}

\begin{figure}[t]
  \centering
  \includegraphics[width=0.7\linewidth]
  {refusal_vs_unsafe.png}
  \caption{Text refusal does not transfer to tool-call safety. For each
  model, verbal refusal rate (no-tool chat) vs.\ unsafe tool-call rate
  (tool-enabled normal) on the same harmful requests.}
  \label{fig:refusal-gap}
\end{figure}

\section{Behavioral Results}

We show the difference between the verbal refusal rate on harmful text prompts with the unsafe tool-call rate on the same tool enabled prompt in Figure~\ref{fig:refusal-gap.}. We observe that all of the five models refuse the majority of harmful requests in chat, and the text refusal rates range from 66\% (Mistral-7B) to 83\% (Qwen3-14B). So basically chat level alignment is broadly similar across model families. However, on the same request, when tools are enabled, every model still makes unsafe tool calls ranging from 10\% (Command-R-7B) to 42\% (Llama-3.1-70B) of the times.
This is the gap between the two bars that appears for every model and is the transfer failure we studu: a high verbal refusal rate does not fully translate into safe action.

In Table~\ref{tab:headline-divergence}, the same gap is quantified over the 384 matched harmful triples per model. The first column shows the no tool text refusal rate and the next two give the unsafe tool-call rate under normal and stress framing, while the final two give the conditioned divergence $\Delta^{\mathrm{cond}}_m$ for each mode. Because we only compute $\Delta^{\mathrm{cond}}_m$ over the triples that the model actually refused in chat, it isolates the transfer failure from the base refusal rate: the text refusal rates mentioned above (0.66 -- 0.82) are pretty high, so these divergence values cannot be ascribed to the fact that one model refuses less often than another. In Normal tool prompts, conditional divergence spans from 0.085 (Command-R-7B) to 0.440 (Llama-3.1-70B). 
Divergence under the stress mode is of comparable magnitude (0.095--0.357) and remains substantial for every model.

These cross-model differences also depend on the frequency with which each model issues tool calls at all, not just how well its refusal transfers: a model that rarely emits a tool call cannot emit an unsafe one. This is explicit in our secondary tool-engagement divergence (any tool call rather than a forbidden one), on chat-refused triples in normal framing, Llama-3.1-70B issues some tool call on roughly 90\% of triples, whereas Command-R-7B does so on only ~14\%. Table~\ref{tab:engagement} reports this any-call divergence, which upper-bounds conditioned divergence; a model's divergence should be read relative to this ceiling rather than in absolute terms.


We show the internal mechanism behind this divergence in Section~\ref{sec:mech}, and show its relationship with model scale and system-prompt framing in Section~\ref{sec:discussion}.

\begin{table}[b]
\centering
\caption{System-prompt framing modulates the transfer gap but does not close it,
shown as conditioned divergence under normal framing by system condition.}
\label{tab:cond}
\begin{tabular}{lccc}
\toprule
Model & Neutral & Safety-reinforced & Tool-encouraging \\
\midrule
Qwen3-14B    & 0.31 & 0.20 & 0.30 \\
Mistral-7B   & 0.12 & 0.14 & 0.14 \\
Command-R-7B & 0.03 & 0.03 & 0.09 \\
\bottomrule
\end{tabular}
\end{table}

\section{Mechanistic Analysis}
\label{sec:mech}

Now we consider why there is a failure and refusal to transfer from text to tool prompts, and whether a single mechanism explains the divergence across model families. Following Section 2.6, we extract the refusal direction for each model and characterize what it does under tool contexts. Table 3 summarizes our headline results.

\begin{table}[b]
\centering
\caption{The refusal mechanism repeats across model families. Suppression is the paired
drop in refusal activation from no-tool to tool-normal, reported as mean $\Delta$ and
$t$-statistic with $n{=}96$. Ablation and addition give refusal rates before and after
the intervention, on harmful and benign prompts respectively. AUC is the area under the
ROC of low refusal activation predicting an unsafe call. Patching is the fraction of
unsafe calls that revert to safe when the projection is restored to its matched no-tool
level.}
\label{tab:mech}
\setlength{\tabcolsep}{4pt}
\begin{tabular}{lccccc}
\toprule
Model & Layer & Suppression ($\Delta$, $t$) & Ablation & Addition & AUC\,/\,Patching \\
\midrule
Qwen3-14B    & 33/40 & $66.8,\ t{=}5.6$ & $54\%\!\to\!34\%$ & $1\%\!\to\!71\%$ & $0.724$\,/\,$18\%$ \\
Mistral-7B   & 26/32 & $1.4,\ t{=}5.1$  & $26\%\!\to\!13\%$ & $0\%\!\to\!80\%$ & $0.689$\,/\,$24\%$ \\
Command-R-7B & 26/32 & $26.4,\ t{=}9.7$ & $70\%\!\to\!2\%$  & $2\%\!\to\!92\%$ & $0.751$\,/\,$62\%$ \\
\bottomrule
\end{tabular}
\end{table}

\paragraph{A single direction causally governs text refusal:}  across every model, we find that ablating our hat sharply decreases the harmful prompt refusal rate and that adding the direction makes the model refuse even benign prompts. In this way, we reproduce the results of \citet{arditi2024refusal} on each architecture. The auto-selected layers, chosen following the method described in Section 2.6, generally  sit in the upper middle part of each network. Furthermore, the addition coefficient scales over two orders of magnitude across the models, ranging from 29 in Mistral to 1,292 in Qwen.

\paragraph{Tool context suppresses the refusal signal:}  for our harmful triples, we find that the mean refusal activation drops significantly from no-tool to tool-enabled modes across all three models (Table 3, with $t = 5.6$ for Qwen, $t = 5.1$ for Mistral, and $t = 9.7$ for Command-R; Figure 1). We find that the text-derived and tool-derived directions remain aligned, revealing that tools attenuate the magnitude of the signal along a shared axis as opposed to rotating it.

\paragraph{The refusal direction serves as a strong predictor but a partial mediator:}  We find that a lower refusal activation reliably predicts an unsafe call in every model, with AUC from 0.69 to 0.75. At the same time, however, when we patch a tool run's projection back to its no-tool level, only a fraction of the unsafe calls revert to being safe: 18 percent for Qwen and 24 percent for Mistral. In this way, we see that transfer failure is not driven solely by refusal suppression alone but also by tool affordance itself. The balance between the two varies by model. For Command-R, for example, the affordance mediates much more strongly than the others at 62 percent, although only on a smaller unsafe set.

\paragraph{Steering is a blunt lever:}  when raising the steering coefficient, the harmful unsafe rate does decrease towards zero, but the benign over-refusal rate also rises accordingly (Figure 2). For example, for Qwen, removing unsafe calls entirely by increasing the steering coefficient costs around 70 percent benign over-refusal. In this case, it would identify any coefficient that removes unsafe behavior without large collateral over-refusal in benign cases. In this way, steering through refusal direction serves as an imprecise safety control.

\begin{figure}[t]
\centering
\begin{minipage}{0.48\textwidth}\centering
  \includegraphics[width=\linewidth]{figures/fig_projection_by_mode.png}
  \caption{Refusal activation by mode (Qwen3-14B). Tool context suppresses the
  projection relative to no-tool chat.}\label{fig:proj}
\end{minipage}\hfill
\begin{minipage}{0.48\textwidth}\centering
  \includegraphics[width=\linewidth]{figures/fig_steering_scaled.png}
  \caption{Steering trade-off (Qwen3-14B). Unsafe calls fall only as benign
  over-refusal rises.}\label{fig:steer}
\end{minipage}
\end{figure}

\section{Discussion}
\label{sec:discussion}


  \paragraph{Text refusal does not predict tool-call safety.} We find that a model's text
  refusal rate is a poor guide to how safely it acts. Llama-3.1-70B refuses 78\%
  of harmful requests in chat but diverges on 44\% of those refused requests once
  tools are available, about five times the divergence of Command-R-7B, which
  refuses at a higher rate (Table~\ref{tab:headline-divergence}). The largest
  model in our study is the least safe in action. Notably, although Llama is among the safest models, it is also the most divergent 

  \paragraph{Affordance, not only pressure, drives the gap.} Although adversarial stress certainly increases divergence for most models, as we see through our two tool-induced prompt phrasings, the effect is comparatively small in comparison to the gap that tool availability alone opens. Ordinary tool prompting Llama is already maximally divergent and does not worsen under stress. We also find that system prompt framing (Table~\ref{tab:cond}) does modulate the gap slightly, but still does not close it. In particular, safety-reinforced prompts help some models (Qwen3, $0.31 \to 0.20$) but not others, and tool-encouraging prompts can widen the gap (Command-R, $0.03 \to 0.09$). Therefore, we see that this failure is a property intrinsic to the tool call setting rather than solely a result of adversarial prompting.


  \paragraph{A shared refusal direction, suppressed rather than replaced.} We find that the single refusal direction that our [Arditi] identifies for chat refusal similarly governs refusal for all of the models we probe: ablations collapse, ablations decrease harmful prompt refusal rates, while additions induce benign over-refusal. Moreover, we find that the directions in the text-derived and tool-derived context are aligned (Table 3). In this way, tool context simply lowers the magnitude of the refusal signal along the same shared axis rather than rotating it. Hence, the divergence reflects a suppression of the existing safety representation rather than a separate tool mode circuit.
  

  \paragraph{Conservative measurement.} By design, our forbidden action checks only fire on arguments that are clearly unsafe. In particular, if a model invents a plausible but incorrect authorization value, it does not count as a violation. Thus, our reported divergence rates should be viewed as lower bounds, and the true gap is likely larger. 

  \paragraph{Implications for deployment.} Our behavioral results show that safety evaluations in text-only settings overestimate the safety of tool-calling agents. Furthermore, since this gap is not fully closed by safety-reinforced prompting but is instead rooted in the fundamental step of action selection, any agents bound for regulated high-stakes environments should be evaluated at the tool-call level in addition to text.

 \section{Future Work}

  \paragraph{Characterizing the affordance pathway.} Restoring the refusal
  direction reverts only part of the unsafe behavior, so a direction-independent
  mechanism tied to tool affordance remains unexplained. Identifying it, whether
  attention to the tool schema, format priors that favor emitting a call, or a
  competing helpfulness feature, is the next mechanistic target.

  \paragraph{Precise interventions.} Single-direction steering trades unsafe calls
  against benign over-refusal too bluntly to deploy. Layer- or token-localized
  interventions, or steering gated on the predicted unsafe-call probability, could
  reduce unsafe actions while preserving benign tool use.

  \paragraph{Harder environments.} The predicates are deterministic single-shot
  checks. Future evaluations should catch fabricated authorization values, extend
  to multi-turn and executed tool calls, and test indirect prompt injection that
  arrives through tool outputs rather than the user turn.

  \paragraph{Refusal monitoring at run time.} The refusal activation predicts
  unsafe calls with AUC up to $0.75$. Even where steering is too blunt, reading the
  projection at the action-selection token gives a runtime signal that an agent is
  about to act unsafely.

  \section{Conclusion}



  We tested whether the safety refusals language models learn in chat transfer to
  their tool-call decisions. Using a 2304-prompt benchmark of matched triples
  across four high-stakes domains, scored by deterministic forbidden-action
  predicates, and a conditioned-divergence metric that separates transfer failure
  from base refusal rates, we find that text refusal does not reliably transfer to
  tool-call safety across model families. Divergence reaches 44\% of refused
  requests, and text-refusal strength is a poor predictor of safety in action. The
  divergence corresponds to graded suppression of the same refusal direction that
  governs chat refusal, attenuated along a shared axis at action selection. This
  suppression predicts unsafe calls but only partly causes them, since tool
  affordance adds an independent pathway. Tool-call safety must be evaluated and
  aligned separately from text safety before LLM agents are deployed in regulated
  environments.

\bibliographystyle{colm2026_conference}
\bibliography{refs}

\vspace{140}
\appendix
\section{Appendix}

\begin{table}[t]
\centering
\caption{Tool-engagement divergence bounds conditioned divergence. For each model we
report conditioned divergence $\Delta^{\mathrm{cond}}_m$ (matched tool run is
\emph{unsafe}) alongside tool-engagement divergence (matched tool run emits \emph{any}
tool call), both conditioned on the chat-refused triples. The gap between them is safe or
benign tool use; the engagement column shows that low divergence can reflect sparse tool
use rather than faithful refusal transfer.}
\label{tab:engagement}
\begin{tabular}{lcccc}
\toprule
& \multicolumn{2}{c}{Normal} & \multicolumn{2}{c}{Stress} \\
\cmidrule(lr){2-3}\cmidrule(lr){4-5}
Model & Divergence & Any-call & Divergence & Any-call \\
\midrule
Llama-3.1-70B-Instruct   & 0.440 & 0.903 & 0.357 & 0.843 \\
Mistral-7B-Instruct-v0.3 & 0.162 & 0.281 & 0.217 & 0.486 \\
Gemma-3-27B-it           & 0.099 & 0.202 & 0.109 & 0.221 \\
Qwen3-14B                & 0.091 & 0.164 & 0.057 & 0.129 \\
c4ai-command-r7b         & 0.085 & 0.142 & 0.095 & 0.206 \\
\bottomrule
\end{tabular}
\end{table}



\end{document}
